\begingroup
\renewcommand{\arraystretch}{1.45}
\Needspace{37\baselineskip}
\InterpolationArticle{39}

Wenn zweitens die Funktion $X$ von der Form
\[
\begin{aligned}
&\beta'\sin x+\beta''\sin 2x+\beta'''\sin 3x+\mathrm{usw.}+\beta^n\sin nx
\end{aligned}
\]
und somit die Anzahl der unbekannten Koeffizienten $=n$ ist, so muss auch die Anzahl der gegebenen Werte,
\SourceLeaf{332}
nach Abzug der überflüssigen, $=n$ sein, damit sie zur vollständigen Bestimmung der Koeffizienten hinreichen. Da nun in diesem Fall diejenigen Werte der Funktion $X$ als überflüssig zu verwerfen sind, die einem Wert von $x$ größer als $180^0$ sowie den Werten $0$ und $180^0$ entsprechen, werden wir für die vier oben unterschiedenen Fälle haben:
\[
\begin{array}{rllcl}
1.&\text{wenn }\mu\text{ gerade},&a=0,&\text{gilt}&n=\frac12\mu-1\\
2.&\text{wenn }\mu\text{ gerade},&a=\frac{180^0}{\mu},&\text{gilt}&n=\frac12\mu\\
3.&\text{wenn }\mu\text{ ungerade},&a=0.&\text{gilt}&n=\frac12\mu-\frac12\\
4.&\text{wenn }\mu\text{ ungerade},&a=\frac{180^0}{\mu},&\text{gilt}&n=\frac12\mu-\frac12
\end{array}
\]
\Needspace{34\baselineskip}
Nun werden ganz ähnlich wie in Artikel 37 in der Funktion
\[
\begin{aligned}
&\gamma+\gamma'\cos x+\gamma''\cos 2x+\text{usw.}+\gamma^m\cos mx\\[6pt]
&{}+\delta'\sin x+\delta''\sin 2x+\text{usw.}+\delta^m\sin mx
\end{aligned}
\]
die nach den Artikeln 20, 22 gewonnen wird, allen $\mu$ Werten genügt und in der $m$ entweder $=\frac12\mu-\frac12$ oder $=\frac12\mu$ ist, die Koeffizienten $\gamma$, $\gamma'$, $\gamma''$ usw. von selbst verschwinden. Denn da in der Reihe $A,B,C,D$ usw. die letzten Terme in umgekehrter Reihenfolge im vorliegenden Fall entweder $=-B,-C,-D$ usw. oder $=-A,-B,-C$ usw. werden, je nachdem, ob $a=0$ oder $=\frac{180^0}{\mu}$ ist, und überdies im ersteren Fall $A=0$ gilt, wird offenbar
\[
\begin{aligned}
&A\cos\lambda a+B\cos\lambda b+C\cos\lambda c+D\cos\lambda d+\text{usw.}
\end{aligned}
\]
für jeden beliebigen Wert von $x$ $=0$ sein. \Needspace{46\baselineskip}
Daher wird die Funktion $X'=$
\[
\begin{aligned}
&\delta'\sin x+\delta''\sin 2x+\delta'''\sin 3x+\text{usw.}+\delta^m\sin mx
\end{aligned}
\]
in der die Koeffizienten $\delta'$, $\delta''$, $\delta'''$ usw. durch die Gleichungen
\[
\begin{aligned}
&\delta'=\frac2\mu(A\sin a+B\sin b+C\sin c+D\sin d+\text{usw.})\\[6pt]
&\delta''=\frac2\mu(A\sin 2a+B\sin 2b+C\sin 2c+D\sin 2d+\text{usw.})\\[6pt]
&\delta'''=\frac2\mu(A\sin 3a+B\sin 3b+C\sin 3c+D\sin 3d+\text{usw.})
\end{aligned}
\]
usw., der letzte aber, wenn $\mu$ gerade und somit $m=\frac12\mu$ ist, durch diese
\[
\begin{aligned}
&\delta^m=\frac1\mu(A\sin ma+B\sin mb+C\sin mc+D\sin md+\text{usw.})
\end{aligned}
\]
bestimmt werden, notwendig mit $X$ identisch sein, was sich ebenso wie in Artikel 37 leicht beweisen lässt. Übrigens steigt die Funktion $X'$ im ersten Fall, wo $a=0$ und $\mu$ gerade ist, tatsächlich nur bis zur Or\SourceLeaf{333}%
dnung $\frac12\mu-1$ an, da
\[
\begin{aligned}
&\delta^m=\frac1\mu(A\sin 0+B\sin 180^0+C\sin 360^0+D\sin 540^0+\text{usw.})\\[6pt]
&=0
\end{aligned}
\]
\Needspace{46\baselineskip}
\InterpolationArticle{40}
Wenn die Funktion $X$ nicht von der $n^{\mathrm{ten}}$ Ordnung ist, wie wir angenommen haben, sondern weiterläuft, so werden die folgenden Gleichungen zeigen, wie der Unterschied zwischen $X'$ und $X$ von den folgenden Koeffizienten abhängt (s. Artikel 21, 23):
\[
\begin{aligned}
&\delta'=\beta'\mp\beta^{\mu-1}\pm\beta^{\mu+1}-\beta^{2\mu-1}+\beta^{2\mu+1}\mp\text{usw.}\\[6pt]
&\delta''=\beta''\mp\beta^{\mu-2}\pm\beta^{\mu+2}-\beta^{2\mu-2}+\beta^{2\mu+2}\mp\text{usw.}\\[6pt]
&\delta'''=\beta'''\mp\beta^{\mu-3}\pm\beta^{\mu+3}-\beta^{2\mu-3}+\beta^{2\mu+3}\mp\text{usw.}
\end{aligned}
\]
und so fort bis zum letzten Koeffizienten $\delta^m$ oder zum vorletzten $\delta^{m-1}$, je nachdem, ob $\mu$ ungerade oder gerade ist; die oberen Zeichen gelten hier, wenn $a=0$ ist, die unteren, wenn $a=\frac{180^0}{\mu}$ ist: Schließlich gilt für den letzten Koeffizienten im Fall (1) $\delta^m=0$, im Fall (2) dagegen
$\delta^m=\beta^m-\beta^{3m}+\beta^{5m}-\beta^{7m}+\text{usw.}$
Alle periodischen Funktionen, die den $\mu$ vorgelegten Werten genügen und unter denen $X'$ die einfachste ist, werden in der Form $X'+Y\sin(\frac12\mu x-\frac12\mu a)$ oder allgemeiner in der Form $X''+Y\sin(\frac12\mu x-\frac12\mu a)$ enthalten sein, wobei $X''$ irgendeine solche Funktion bezeichnet; diese Formel geht für die Fälle 1 und 3 auf $X''+Y\sin\frac12\mu x$, für die Fälle 2 und 4 dagegen auf $X''+Y\cos\frac12\mu x$ zurück. Aber $Y$ muss, da wir keine anderen Funktionen betrachten als solche, die $X$ ähnlich sind, d. h. nur aus Sinusgliedern bestehen, notwendig sein:
\[
\begin{array}{ll}
\text{im Fall 1 von der Form}&g+g'\cos x+g''\cos 2x+\text{usw.}\\
\text{im Fall 2 von der Form}&g\sin x+g'\sin 2x+g''\sin 3x+\text{usw.}\\
\text{im Fall 3 von der Form}&g\cos\frac12x+g'\cos\frac32x+g''\cos\frac52x+\text{usw.}\\
\text{im Fall 4 von der Form}&g\sin\frac12x+g'\sin\frac32x+g''\sin\frac52x+\text{usw.}
\end{array}
\]
\Needspace{26\baselineskip}
Die Funktion $Y$ wird sich hier so bestimmen lassen, dass eine auf die Ordnung $\frac12\mu-1$, $\frac12\mu$, $\frac12\mu-\frac12$, $\frac12\mu-\frac12$ herabgesetzte Funktion hervorgeht, die notwendig mit $X'$ identisch sein wird. Für die Herabsetzung der Funktion $X''$ auf $X'$ gilt folgende allgemeine Regel:

Jeder Term in $X''$ von der Form $L\sin\lambda x=L\sin(k\mu\pm\lambda')x$ werde entweder in $\pm L\sin\lambda'x$ verwandelt (sooft $a=0$ oder $k$ gerade ist) oder in $\mp L\sin\lambda'x$ (sooft weder $a=0$ noch $k$ gerade ist, d. h. sooft zugleich $a=\frac{180^0}{\mu}$ und $k$ ungerade ist): Schließlich werde, sooft im
\SourceLeaf{334}
ersten Fall, d. h. für \(a=0\) und einen geraden Wert von \(\mu\), \(\lambda'=\frac12\mu\) wird, der Term \(L\sin\lambda x\) ganz vernichtet.

\endgroup
